%% kapa: exam=1 round=21 year=2010 period=1 form=A subject=economics subject_ko=경제원론 seq=21 number=61
%% source: Q-Net 감정평가사 기출문제 (cst00309) artlSeq=5024181 file=extracted/1008/5024181/2010년 제21회 감정평가사 1차시험 기출문제/1교시.hwp
%% conversion: kordoc 4.24.51 -> tools/gen1.py (auto); review=proofread-sonnet; flags=none
\begin{question}{61}
완전경쟁시장에 참여하는 모든 기업의 비용함수가 동일하며 평균비용곡선이 U자형이다. 다음 설명 중 옳지 않은 것은?
\begin{choices}
  \item 기업은 가격수용자로서 행동한다.
  \item 단기에 경제적 이윤이 발생할 수 있다.
  \item 기업의 진입ㆍ퇴출이 자유로운 장기에는 경제적 이윤은 0이다.
  \item 장기균형에서 가격은 한계비용과 같다.
  \item 장기균형에서 한계비용은 평균비용보다 높다.
\end{choices}
\end{question}
\answer{5}
